\documentclass[fleqn,10pt]{wlscirep}

\usepackage[utf8]{inputenc}
\usepackage[T1]{fontenc}
\usepackage{booktabs}
\usepackage{multirow}
\usepackage{amsmath}
\usepackage{amssymb}
\usepackage{graphicx}
\usepackage{xcolor}
\usepackage{hyperref}
\usepackage{array}
\usepackage{caption}
\usepackage{subcaption}
\usepackage{placeins}

\renewcommand{\topfraction}{0.85}
\renewcommand{\bottomfraction}{0.85}
\renewcommand{\textfraction}{0.10}
\renewcommand{\floatpagefraction}{0.80}
\setcounter{topnumber}{4}
\setcounter{bottomnumber}{4}
\setcounter{totalnumber}{6}

% ── Colour palette ──────────────────────────────────────────────────────────
\definecolor{agentblue}{HTML}{185FA5}
\definecolor{agentteal}{HTML}{1D9E75}
\definecolor{agentcoral}{HTML}{E8593C}

% ── Convenience macros ──────────────────────────────────────────────────────
\newcommand{\env}{\textsc{LakeGym~v5}}
\newcommand{\attentive}{\texttt{AttentivePPO}}
\newcommand{\mlpppo}{\texttt{MLP-PPO}}
\newcommand{\ddqn}{\texttt{DDQN}}
\newcommand{\noop}{\texttt{NOOP}}
\newcommand{\cthirtytwo}{\texttt{COMPACT\_32KB}}
\newcommand{\csixtyfour}{\texttt{COMPACT\_64KB}}
\newcommand{\conehundredtwentyeight}{\texttt{COMPACT\_128KB}}
\newcommand{\phour}{\texttt{PARTITION\_HOUR}}
\newcommand{\pregion}{\texttt{PARTITION\_REGION}}
\newcommand{\ptype}{\texttt{PARTITION\_EVENT\_TYPE}}
\newcommand{\removep}{\texttt{REMOVE\_PARTITION}}

% ── Title & authors ─────────────────────────────────────────────────────────
\title{Workload-Aware Data Lakehouse Maintenance Using Hierarchical Deep
Reinforcement Learning}

\author[1,*]{Omer Karim}
\author[1,*]{Omar Essameldin}
\author[1]{Omar Shrief}
\author[1,2]{Ahmed B. Zaky}

\affil[1]{Department of Computer Science and Information Technology,
  Egypt-Japan University of Science and Technology, Alexandria, Egypt}
\affil[2]{Department of Electrical Engineering, Faculty of Engineering
  (Shoubra), Benha University, Banha 13511, Qalubiya, Egypt}

\affil[*]{These authors contributed equally to this work.}

\affil[$\dagger$]{\texttt{omer.karim@ejust.edu.eg},
\texttt{omar.isameldin@ejust.edu.eg},
\texttt{omar.shrief@ejust.edu.eg},
\texttt{ahmed.zaky@ejust.edu.eg}}
\affil[$\ddagger$]{Correspondence should be addressed to O.E. 
(\texttt{omar.isameldin@ejust.edu.eg})}
% ── Abstract ────────────────────────────────────────────────────────────────
\begin{abstract}
Modern data lakehouses built on open table formats such as Apache Iceberg
accumulate small files and suboptimal partition layouts as data is
continuously ingested, degrading query performance over time.
Existing maintenance strategies rely on hand-crafted, workload-agnostic
rules that cannot adapt to shifting query distributions.
We present a hierarchical multi-agent deep reinforcement learning (RL)
framework that jointly optimises compaction and partitioning decisions in
real time.
A rule-based meta-controller evaluates per-step urgency scores and
delegates to two specialist RL agents---a compaction agent and a
partition agent---each trained with a separate reward signal aligned to
its operational objective.
We evaluate three neural architectures within this framework: an
attention-based actor-critic (\attentive{}), a feedforward actor-critic
(\mlpppo{}), and a dueling double DQN (\ddqn{}).
Experiments on \env{}, a simulation environment backed by a real Iceberg
table under a structured 1,000-step workload with five distinct query
profiles, show that all three hierarchical RL agents outperform 17
heuristic and workload-aware baselines.
\attentive{} achieves the highest mean global reward (0.2201), a
$+$4.3\% improvement over the best heuristic and $+$29.5\% over a
compaction-only ablation.
\ddqn{} achieves the lowest estimated cloud cost (\$8.53 per episode,
$-$82.9\% versus no-maintenance), with query latency (323.4\,ms) nearly
identical to \attentive{} (323.6\,ms) across the five-episode average.
Ablation experiments confirm that compaction is a prerequisite for
effective partitioning: a partition-only agent produces catastrophic
file proliferation (1,245 mean files, 5,307\,ms latency).
Adaptive online weight updates further improve \attentive{} by
$+$20.3\% across five evaluation episodes, closing 146.9\% of the
initial gap to its frozen ceiling, while destabilising \ddqn{}, exposing
an important architecture-specific constraint on on-policy versus
off-policy online adaptation.
\end{abstract}

% ════════════════════════════════════════════════════════════════════════════
\begin{document}
\flushbottom
\maketitle
\thispagestyle{empty}

% ════════════════════════════════════════════════════════════════════════════
\section*{Introduction}
\label{sec:intro}
% ════════════════════════════════════════════════════════════════════════════

Data lakehouses combine the scalable storage of data lakes with the
transactional semantics of data warehouses through open table formats
such as Apache Iceberg\textsuperscript{1}, Delta
Lake\textsuperscript{2}, and Apache Hudi\textsuperscript{3}.
These formats store data as immutable Parquet files and maintain
metadata-level transaction logs, enabling ACID semantics on object
stores.
However, continuous ingestion generates large numbers of small files,
and the initial partition layout quickly becomes misaligned with the
actual query distribution as workloads evolve, increasing query scan
costs substantially.

State-of-practice maintenance is either manual or based on periodic
compaction schedules that are decoupled from query patterns.
Open-source auto-compaction heuristics\textsuperscript{1} trigger on
file-count thresholds but do not consider partition alignment.
Workload-aware approaches such as learned index
structures\textsuperscript{4} and automatic physical design
tools\textsuperscript{5} address layout optimisation offline and are
not designed for the continuous-ingestion regime of a lakehouse.
Reinforcement learning has been applied to query
optimisation\textsuperscript{6,7} and database
tuning\textsuperscript{8,9}, but not to joint compaction and
partitioning maintenance.

We argue that effective lakehouse maintenance requires \emph{joint},
\emph{online}, and \emph{workload-aware} decisions over both
compaction and partitioning.
This paper makes the following contributions.

\begin{itemize}
\item \textbf{Framework.} A hierarchical multi-agent RL framework with
  a rule-based meta-controller that separates the when-to-act decision
  from the what-to-do decision, improving interpretability and training
  stability.

\item \textbf{Environment.} \env{}, a simulation environment backed by
  a real Apache Iceberg table on Apache Spark with a structured,
  phase-shifting workload and a three-component reward function.

\item \textbf{Architecture study.} A systematic comparison of three
  neural architectures (\attentive{}, \mlpppo{}, \ddqn{}) in both
  frozen and adaptive evaluation modes across 22 policies.

\item \textbf{Ablation.} Controlled experiments confirming that joint
  compaction and partitioning is non-trivially superior to either
  sub-agent acting alone, with $+$26--33\% reward gain over
  compaction-only and $+$221--228\% over partition-only baselines.
\end{itemize}

\FloatBarrier

% ════════════════════════════════════════════════════════════════════════════
\section*{Results}
\label{sec:results}
% ════════════════════════════════════════════════════════════════════════════

\subsection*{Overall policy ranking}

Table~\ref{tab:ranking} presents the mean global reward, mean latency,
mean file count, and mean pruning ratio for all 22 evaluated policies
on the 1,000-step standard workload.
All three hierarchical RL agents occupy the top three positions.
\attentive{} achieves the highest mean global reward (0.2201), followed
by \ddqn{} (0.2196) and \mlpppo{} (0.2124).
The best workload-aware heuristic,
\texttt{WorkloadAwareThreshold\_v5}, ranks fourth (0.2111),
representing a $+$4.3\% improvement for the best RL agent.
Relative to the best pure-compaction heuristic
(\texttt{AlwaysCompact\_C128}, reward $= 0.1869$), \attentive{}
provides a $+$17.8\% gain.
No-maintenance is the worst policy at $-0.1492$, confirming that
inaction is severely penalised under continuous ingestion.

\begin{table}[!htbp]
\centering
\setlength{\tabcolsep}{5pt}
\begin{tabular}{clcccc}
\toprule
\textbf{Rank} & \textbf{Policy} & \textbf{Reward} &
  \textbf{Latency (ms)} & \textbf{Files} & \textbf{Pruning} \\
\midrule
1  & \textbf{\attentive{}}           & \textbf{0.2201} & 323.6           & 54.9   & 0.191 \\
2  & \textbf{\ddqn{}}                & 0.2196          & \textbf{323.4}  & 58.8   & 0.197 \\
3  & \textbf{\mlpppo{}}              & 0.2124          & 387.8           & 62.7   & \textbf{0.203} \\
4  & \texttt{WorkloadAwareThreshold} & 0.2111          & 439.9           & 93.4   & 0.285 \\
5  & \texttt{AlwaysCompact\_C128}    & 0.1869          & 344.3           & 43.5   & 0.000 \\
6  & \texttt{Random\_Compact}        & 0.1809          & 460.8           & 52.1   & 0.000 \\
7  & \texttt{Threshold10\_C128}      & 0.1734          & 429.7           & 46.8   & 0.000 \\
8  & \texttt{Periodic10\_C128}       & 0.1734          & 399.1           & 47.6   & 0.000 \\
9  & \texttt{CompactOnly} (ablation) & 0.1699          & 344.1           & 48.9   & 0.000 \\
10 & \texttt{PartitionExploration}   & 0.1668          & 437.2           & 90.8   & 0.138 \\
11 & \texttt{AdaptiveCompaction}     & 0.1411          & 609.5           & 71.4   & 0.000 \\
12 & \texttt{AlwaysCompact (C64)}    & 0.1393          & 554.0           & 88.0   & 0.000 \\
\midrule
21 & \texttt{No\_Maintenance}        & $-$0.1492       & 4,298.1         & 873.7  & 0.000 \\
22 & \texttt{PartitionOnly} (ablation)& $-$0.1765      & 5,306.9         & 1,245.0 & 0.110 \\
\bottomrule
\end{tabular}
\caption{\label{tab:ranking}Full policy ranking. Mean global reward, mean latency (ms),
  mean file count, and mean pruning ratio for all 22 evaluated policies
  on the 1,000-step v5 standard workload, sorted by descending reward.
  RL agents are shown in bold. Top-12 and bottom-2 policies shown;
  the complete table is available as supplementary data.}
\end{table}

Figure~\ref{fig:cumreward} shows the cumulative global reward trajectory
over 1,000 steps. The RL ensemble diverges sharply upward from all
baselines during the burst-ingestion window (steps 301--500), where
the ingestion rate reaches 225\,rows/step, demonstrating the agents'
ability to recognise and respond to increased storage pressure in real
time.

\begin{figure}[!htbp]
\centering
\includegraphics[width=0.92\linewidth]{charts/fig_cumreward}
\caption{\label{fig:cumreward}\textbf{Cumulative reward race over 1,000 steps.}
  Solid lines show the three hierarchical RL agents
  (\attentive{} = blue, \mlpppo{} = green dashed, \ddqn{} = red/orange dashed).
  Dashed lines show selected heuristic baselines
  (no-maintenance = light red, always-compact = gray,
  adaptive-compact = yellow-green, threshold10-C64 = yellow).
  The light-blue shaded band spans the RL agent ensemble range.
  All RL agents diverge sharply upward from heuristic baselines across
  the full 1,000-step episode.
  Final cumulative rewards: \attentive{} = 220.09,
  \ddqn{} = 219.60, \mlpppo{} = 212.37 versus best heuristic = 141.05.}
\end{figure}

Figure~\ref{fig:heuristic_baseline} provides the heuristic-only
comparison across three metrics.

\begin{figure}[!htbp]
\centering
\includegraphics[width=0.92\linewidth]{charts/fig_heuristic_baseline}
\caption{\label{fig:heuristic_baseline}\textbf{Heuristic baseline comparison on the 1,000-step v5 workload.}
  Horizontal bar charts showing mean global reward (left), mean latency
  (middle), and mean file count (right) for all heuristic policies,
  ranked by descending reward. \texttt{AlwaysCompact\_C128} achieves the
  best reward (0.1869) and lowest latency (344\,ms) among all heuristics.
  \texttt{No\_Maintenance} is the worst policy on all three metrics
  (reward $= -0.1492$, latency $= 4{,}298$\,ms, files $= 874$).}
\end{figure}

\FloatBarrier

\subsection*{Three-way RL agent comparison}

Table~\ref{tab:threeway} compares the three RL architectures on four
metrics. \attentive{} achieves the best mean global reward (0.2201) and
lowest mean file count (54.9), while \ddqn{} achieves the lowest mean
latency (323.4\,ms), though the margin over \attentive{} (323.6\,ms) is
negligible across the five-episode average. \mlpppo{} achieves the
highest mean pruning ratio (0.2028) but ranks last on reward and latency.
Step-level global reward statistics (mean $\pm$ s.d.) are
$0.2201 \pm 0.0646$ for \attentive{}, $0.2124 \pm 0.0643$ for
\mlpppo{}, and $0.2196 \pm 0.0650$ for \ddqn{}; the tighter standard
deviations relative to single-episode figures represents the smoothing effect
of the five-episode averaging.
The pairwise reward gap between \attentive{} and \ddqn{} is $-0.0005$
($-$0.2\%), and between \attentive{} and \mlpppo{} is $-0.0077$
($-$3.5\%). These modest gaps, together with the distinct per-metric
strengths of each agent (Section~\ref{sec:crossarch}), suggest that no
single architecture dominates across all operational objectives.

\begin{table}[!htbp]
\centering
\begin{tabular}{lcccc}
\toprule
\textbf{Agent} & \textbf{Reward} & \textbf{Latency (ms)} &
  \textbf{Files} & \textbf{Pruning} \\
\midrule
\attentive{} & $0.2201 \pm 0.0646$ & $323.6 \pm 143.3$ & $54.9 \pm 28.0$ & $0.191 \pm 0.128$ \\
\ddqn{}      & $0.2196 \pm 0.0650$ & $323.4 \pm 142.9$ & $58.8 \pm 30.2$ & $0.197 \pm 0.137$ \\
\mlpppo{}    & $0.2124 \pm 0.0643$ & $387.8 \pm 171.5$ & $62.7 \pm 30.9$ & $0.203 \pm 0.148$ \\
\bottomrule
\end{tabular}
\caption{\label{tab:threeway}Three-way RL agent comparison. Mean global reward, mean latency
  (ms), mean file count, and mean pruning ratio for the three frozen RL
  agents on the 1,000-step v5 standard workload, reported as step-level
  mean $\pm$ s.d.\ averaged over five independent episodes ($n = 5{,}000$
  total steps per agent).}
\end{table}

Figure~\ref{fig:rlcomp} shows the four-metric comparison for all three
agents, and Figure~\ref{fig:rl_dynamics} shows the step-by-step
dynamics over the full 1,000-step standard workload.

\begin{figure}[!htbp]
\centering
\includegraphics[width=0.82\linewidth]{charts/fig_rlcomp}
\caption{\label{fig:rlcomp}\textbf{Three-way RL agent comparison on the 1,000-step
  standard workload.}
  \textbf{(a)}~Mean global reward.
  \textbf{(b)}~Mean query latency (ms).
  \textbf{(c)}~Mean file count.
  \textbf{(d)}~Mean average pruning ratio.
  Bars use the agent colour convention: blue = \attentive{},
  teal = \mlpppo{}, coral = \ddqn{}.
  Value annotations are shown above each bar.
  \attentive{} leads on reward and file management; \ddqn{} leads
  on latency; \mlpppo{} leads on pruning ratio.}
\end{figure}

\begin{figure}[!htbp]
\centering
\includegraphics[width=0.92\linewidth]{charts/fig_rl_dynamics}
\caption{\label{fig:rl_dynamics}\textbf{Frozen RL agents --- step-by-step dynamics over the
  1,000-step standard workload.}
  Four panels show global reward (top-left), query latency (top-right),
  file count (bottom-left), and average pruning ratio (bottom-right).
  Lines use the agent colour convention:
  \attentive{} = blue, \mlpppo{} = teal, \ddqn{} = coral. All three
  agents maintain tightly clustered reward and latency profiles
  throughout the workload.}
\end{figure}

\FloatBarrier

\subsection*{Reward decomposition}

Table~\ref{tab:rewarddecomp} decomposes mean reward into its three
components. All three RL agents are the only policies achieving a
positive mean partition reward (0.051--0.068). All heuristic compaction
policies produce negative partition rewards (e.g., \texttt{AlwaysCompact}:
$-0.050$) because they never change the partition layout. The ablation
policies confirm structural roles: \texttt{CompactOnly} achieves the
highest compact reward (0.1855) but zero partition reward, while
\texttt{PartitionOnly} achieves marginal partition reward ($-0.018$)
but catastrophically negative compact reward ($-0.140$) owing to
uncontrolled file proliferation. \ddqn{} achieves the highest partition
reward of any agent (0.0675), suggesting its Q-value formulation may
better exploit long-horizon pruning benefits.

\begin{table}[!htbp]
\centering
\begin{tabular}{lccc}
\toprule
\textbf{Policy} & \textbf{Compact reward} &
  \textbf{Partition reward} & \textbf{Global reward} \\
\midrule
No maintenance           & $-0.0851$           & $\phantom{+}0.0000$ & $-0.1492$ \\
Always compact           & $\phantom{+}0.3608$ & $-0.0500$           & $\phantom{+}0.1393$ \\
Adaptive compact         & $\phantom{+}0.1417$ & $-0.0428$           & $\phantom{+}0.1411$ \\
Compact only (ablation)  & $\phantom{+}0.1855$ & $-0.0085$           & $\phantom{+}0.1699$ \\
Partition only (ablation)& $-0.1397$           & $-0.0179$           & $-0.1765$ \\
\attentive{}             & $\phantom{+}0.1581$ & $+0.0535$           & $\mathbf{0.2201}$ \\
\mlpppo{}                & $\phantom{+}0.1840$ & $+0.0510$           & $0.2124$ \\
\ddqn{}                  & $\phantom{+}0.1659$ & $\mathbf{+0.0675}$  & $0.2196$ \\
\bottomrule
\end{tabular}
\caption{\label{tab:rewarddecomp}Reward decomposition. Mean compact reward, partition reward,
  and global reward for key policies. Only the three RL agents achieve
  positive mean partition reward.}
\end{table}

\begin{figure}[!htbp]
\centering
\includegraphics[width=0.92\linewidth]{charts/fig_reward_decomp}
\caption{\label{fig:reward_decomp}\textbf{Reward decomposition across key policies.}
  Bar charts showing mean compact reward (left), mean partition reward
  (centre), and mean global reward (right) for all key policies.
  Only the three RL agents (\attentive{}, \mlpppo{}, \ddqn{}) achieve
  positive partition reward; all compaction heuristics yield negative
  partition reward, and \texttt{No\_Maintenance} yields zero (no
  partition changes are ever made).}
\end{figure}

\FloatBarrier

\subsection*{Compaction target size sensitivity}

Table~\ref{tab:csize} shows the effect of compaction target size on
the always-compact baseline. C32\,KB is counter-productive (1,283.0\,ms
latency, 179.1 files) because small target files amplify the number of
compaction operations without reducing overall file count. C128\,KB
produces the best result (344.3\,ms, 43.5 files, reward 0.1869),
motivating the use of \conehundredtwentyeight{} under moderate-to-high
ingestion loads. The trained RL agents internalised this: \attentive{}
selects \conehundredtwentyeight{} in 10.8\% of steps versus \csixtyfour{}
in 5.0\% and \cthirtytwo{} in 2.2\%.

\begin{table}[!htbp]
\centering
\begin{tabular}{lccc}
\toprule
\textbf{Target size} & \textbf{Reward} &
  \textbf{Latency (ms)} & \textbf{Files} \\
\midrule
32\,KB  & 0.0393 & 1,283.0 & 179.1 \\
64\,KB  & 0.1393 & 554.0   & 88.0  \\
128\,KB & \textbf{0.1869} & \textbf{344.3} & \textbf{43.5} \\
\bottomrule
\end{tabular}
\caption{\label{tab:csize}Compaction target size sensitivity. Mean global reward, mean
  query latency (ms), and mean file count for the three always-compact
  variants ($n = 1{,}000$ steps each).}
\end{table}

\begin{figure}[!htbp]
\centering
\includegraphics[width=0.85\linewidth]{charts/fig_csize}
\caption{\label{fig:csize}\textbf{Compaction target size comparison --- C32 vs C64 vs C128.}
  Rolling 20-step mean global reward (left), query latency (centre),
  and file count (right) for the three always-compact variants over
  1,000 steps. C128 (blue) achieves the lowest latency and file count;
  C32 (red) produces severe file proliferation and high latency.}
\end{figure}

\FloatBarrier

\subsection*{Action distribution}

Table~\ref{tab:actiondist} summarises action distributions of the three
frozen agents, pooled across five independent 1,000-step episodes
(5,000 steps per agent). All three remain predominantly idle (\noop{}:
76.2--79.3\% of steps), confirming that selective intervention rather
than compulsive action is a robust property that holds across runs.
Total compaction rates are similar across agents (16.6--18.2\%) while
partition rates are lower (4.1--5.6\%). \attentive{} still issues the
most partition actions (5.6\% per step on average), consistent with its
highest global reward. \ddqn{} now shows the most evenly spread
compaction strategy across all three target sizes, while \mlpppo{}
continues to concentrate almost exclusively on \csixtyfour{} (16.7\%).
Figure~\ref{fig:actiondist} shows the per-agent action share plots.

\begin{table}[!htbp]
\centering
\begin{tabular}{lcccccc}
\toprule
\textbf{Action} &
  \multicolumn{2}{c}{\textbf{\attentive{}}} &
  \multicolumn{2}{c}{\textbf{\mlpppo{}}} &
  \multicolumn{2}{c}{\textbf{\ddqn{}}} \\
  & $n$/ep & \% & $n$/ep & \% & $n$/ep & \% \\
\midrule
NOOP                            & 762.2 & 76.2 & 775.8 & 77.6 & 793.4 & 79.3 \\
\texttt{COMPACT\_128KB}         & 108.2 & 10.8 &  14.2 &  1.4 &  56.4 &  5.6 \\
\texttt{COMPACT\_64KB}          &  46.2 &  4.6 & 167.0 & 16.7 &  49.8 &  5.0 \\
\texttt{COMPACT\_32KB}          &  27.0 &  2.7 &   0.8 &  0.1 &  59.4 &  5.9 \\
\textit{All compaction}         & 181.4 & 18.1 & 182.0 & 18.2 & 165.6 & 16.6 \\
\addlinespace
\texttt{PARTITION\_REGION}      &  34.0 &  3.4 &   8.0 &  0.8 &  10.0 &  1.0 \\
\texttt{PARTITION\_EVENT\_TYPE} &   6.4 &  0.6 &  11.0 &  1.1 &  11.2 &  1.1 \\
\texttt{PARTITION\_HOUR}        &   6.8 &  0.7 &  10.8 &  1.1 &  11.0 &  1.1 \\
\texttt{REMOVE\_PARTITION}      &   9.2 &  0.9 &  12.4 &  1.2 &   8.8 &  0.9 \\
\textit{All partition}          &  56.4 &  5.6 &  42.2 &  4.2 &  41.0 &  4.1 \\
\bottomrule
\end{tabular}
\caption{\label{tab:actiondist}Action distribution of the three frozen RL agents pooled
  over five independent 1,000-step episodes (5,000 steps total per agent).
  Values show per-episode average count ($n$/ep) and percentage share.
  All agents are predominantly idle (NOOP $> 76\%$).}
\end{table}

\begin{figure}[!htbp]
\centering
\begin{subfigure}[b]{0.32\textwidth}
  \includegraphics[width=\linewidth]{charts/fig_action_attentive}
  \caption{\attentive{}}
\end{subfigure}
\hfill
\begin{subfigure}[b]{0.32\textwidth}
  \includegraphics[width=\linewidth]{charts/fig_action_mlpppo}
  \caption{\mlpppo{}}
\end{subfigure}
\hfill
\begin{subfigure}[b]{0.32\textwidth}
  \includegraphics[width=\linewidth]{charts/fig_action_ddqn}
  \caption{\ddqn{}}
\end{subfigure}
\caption{\label{fig:actiondist}\textbf{Action distributions of the three frozen RL agents
  (5 episodes, 5,000 steps each).}
  Horizontal bar charts showing the per-episode average share of each
  action type. All agents are predominantly idle (NOOP $> 76\%$).
  \attentive{} issues the most partition actions (5.6\%); \mlpppo{}
  concentrates almost entirely on \csixtyfour{} (16.7\%); \ddqn{}
  spreads compaction evenly across all three target sizes.}
\end{figure}

\FloatBarrier

\subsection*{Frozen agent consistency}
\label{sec:consistency}

To assess run-to-run reliability, each frozen agent was evaluated over
five independent 500-step episodes on the compressed evaluation workload.
Table~\ref{tab:consistency} reports the coefficient of variation (CV)
across episodes for each agent and metric.
\attentive{} is the most consistent agent: reward CV $= 5.86\%$,
latency CV $= 2.33\%$, file count CV $= 2.40\%$.
\mlpppo{} now shows moderate variance on reward (CV $= 6.99\%$) and,
with the corrected 500-step data, a much tighter latency spread
(CV $= 3.24\%$, mean 258.4\,ms) --- the earlier single-episode anomaly
at 428.8\,ms no longer inflates the variance figure.
\ddqn{} occupies an intermediate position (reward CV $= 6.38\%$,
latency CV $= 6.73\%$).
Pruning ratio remains the most variable metric across all agents
(CV $= 18$--25\%), reflecting the inherent stochasticity of
workload-phase pruning opportunities.
Figure~\ref{fig:frozenconsist} shows the per-episode trajectories.

\begin{table}[!htbp]
\centering
\begin{tabular}{lcccccc}
\toprule
& \multicolumn{2}{c}{\textbf{\attentive{}}} &
  \multicolumn{2}{c}{\textbf{\mlpppo{}}} &
  \multicolumn{2}{c}{\textbf{\ddqn{}}} \\
\textbf{Metric} & Mean & CV & Mean & CV & Mean & CV \\
\midrule
Global reward & 0.2182 & 5.86\%  & 0.1991 &  6.99\% & 0.2065 &  6.38\% \\
Latency (ms)  & 190.3  & 2.33\%  & 258.4  &  3.24\% & 246.2  &  6.73\% \\
File count    & 31.8   & 2.40\%  & 37.2   &  2.88\% & 35.2   &  3.28\% \\
Pruning ratio & 0.214  & 18.26\% & 0.191  & 24.80\% & 0.193  & 24.70\% \\
\bottomrule
\end{tabular}
\caption{\label{tab:consistency}Frozen agent consistency across five evaluation episodes
  ($n = 500$ steps each). CV = coefficient of variation
  (s.d.$/$ mean $\times 100\%$); lower values indicate higher consistency.}
\end{table}

\begin{figure}[!htbp]
\centering
\includegraphics[width=0.92\linewidth]{charts/fig_frozen_consistency}
\caption{\label{fig:frozenconsist}\textbf{Frozen agents --- mean metric per evaluation episode
  (500-step workload, 5 episodes).}
  Mean global reward, mean query latency, mean file count, and mean
  pruning ratio per episode for the three frozen agents.
  Lines use the agent colour convention: \attentive{} = blue,
  \mlpppo{} = teal, \ddqn{} = coral.
  \attentive{} shows the tightest clustering on reward (CV = 5.86\%),
  latency (CV = 2.33\%), and file count (CV = 2.40\%).}
\end{figure}

\FloatBarrier

\subsection*{Adaptive online learning}
\label{sec:adaptive}

Table~\ref{tab:adaptive} reports the episode-over-episode reward
progression for the adaptive variants of all three agents.
\attentive{} shows a strong upward trend from 0.1934 (episode~1) to 0.2326
(episode~5), a total gain of $+0.0393$ ($+$20.3\%), with a brief dip at
episode~4 (0.2149) before recovering.
By episode~5, the adaptive agent \emph{exceeds} its own frozen ceiling
(0.2201) by 0.0125, closing 146.9\% of the initial gap --- a stronger
result than the single-episode comparison suggested, since the averaged
ceiling sits somewhat lower and the adaptive trajectory is unchanged.
\mlpppo{} shows modest improvement ($+$4.9\%), closing 75.1\% of its
gap to the updated frozen ceiling (0.2124). In contrast, \ddqn{} degrades
from 0.2181 to 0.1696 ($-$22.2\%), peaking at episode~3 (0.2298) before
collapsing---a hallmark of replay-buffer instability in off-policy online
adaptation.
Inspection of \ddqn{}'s adaptive action distribution reveals increasing
compaction frequency (from 16.4\% to 20.8\%) concurrent with reward
degradation, suggesting the replay buffer reinforces over-compact
behaviour once encountered. Figure~\ref{fig:adaptive} shows the
reward convergence curves with frozen-ceiling reference lines.

\begin{table}[!htbp]
\centering
\begin{tabular}{lcccccccc}
\toprule
\textbf{Agent} & \textbf{Ep 1} & \textbf{Ep 2} & \textbf{Ep 3} &
  \textbf{Ep 4} & \textbf{Ep 5} & \textbf{Gain} &
  \textbf{Ceiling} & \textbf{\% gap closed} \\
\midrule
\attentive{} & 0.1934 & 0.2017 & 0.2226 & 0.2149 & 0.2326 &
  $+$20.3\% & 0.2201 & 146.9\% \\
\mlpppo{}    & 0.1994 & 0.2018 & 0.1818 & 0.1939 & 0.2091 &
  $+$4.9\%  & 0.2124 & 75.1\%  \\
\ddqn{}      & 0.2181 & 0.2133 & 0.2298 & 0.2047 & 0.1696 &
  $-$22.2\% & 0.2196 & $-$3,207\% \\
\bottomrule
\end{tabular}
\caption{\label{tab:adaptive}Adaptive online learning. Mean global reward per episode
  for the three adaptive agents ($n = 500$ steps each). ``Gain'' reports
  the relative change from episode~1 to episode~5. ``\% gap closed''
  measures progress toward the five-episode-averaged frozen ceiling;
  values above 100\% indicate the adaptive agent exceeded its frozen ceiling.}
\end{table}

\begin{figure}[!htbp]
\centering
\includegraphics[width=0.92\linewidth]{charts/fig_adaptive_episodes}
\caption{\label{fig:adaptive}\textbf{Adaptive agents --- mean metric per episode (online
  learning, 500-step workload).}
  Mean global reward, mean query latency, mean file count, and mean
  pruning ratio per episode for the three adaptive agents.
  \attentive{} (blue) shows a strong upward trend overall, rising from
  0.1934 (episode~1) to 0.2326 (episode~5) with a brief dip at
  episode~4, and ultimately exceeds its frozen ceiling ($+$20.3\%).
  \mlpppo{} (teal) shows modest improvement ($+$4.9\%).
  \ddqn{} (coral) peaks at episode~3 then degrades sharply to $-$22.2\%,
  indicating replay-buffer instability under online adaptation.}
\end{figure}

\FloatBarrier

\subsection*{Frozen versus adaptive comparison}

Table~\ref{tab:fva} shows the per-episode reward comparison between
frozen and adaptive modes for each architecture.
For \attentive{}, the adaptive agent underperforms frozen in episodes
1--3 (by $-$11.9\% to $-$4.9\%) but surpasses it in episodes 4--5
($+$9.4\% and $+$2.5\%), indicating a necessary warm-up period before
online learning converges.
For \mlpppo{}, the adaptive agent outperforms frozen in episode~1
($+$1.1\%) but underperforms in 3 of the remaining 4 episodes,
suggesting the feedforward architecture still learns slowly from
online gradient updates overall.
For \ddqn{}, the adaptive agent outperforms frozen in episodes 1--3 but
collapses in episode~5 ($-$18.9\%), confirming that off-policy
replay-buffer adaptation is unsuitable for the non-stationary evaluation
environment. These results collectively indicate that \textbf{frozen
weights are preferred for deployment}, and online adaptation should be
restricted to on-policy architectures such as \attentive{}.

\begin{table}[!htbp]
\centering
\setlength{\tabcolsep}{4pt}
\begin{tabular}{llccccc}
\toprule
\textbf{Agent} & \textbf{Mode} & \textbf{Ep 1} & \textbf{Ep 2} &
  \textbf{Ep 3} & \textbf{Ep 4} & \textbf{Ep 5} \\
\midrule
\multirow{3}{*}{\attentive{}}
  & Frozen     & 0.2194     & 0.2144    & 0.2340     & 0.1964     & 0.2270 \\
  & Adaptive   & 0.1934     & 0.2017    & 0.2226     & 0.2149     & 0.2326 \\
  & $\Delta\%$ & $-$11.9    & $-$5.9    & $-$4.9     & $+$9.4     & $+$2.5 \\
\midrule
\multirow{3}{*}{\mlpppo{}}
  & Frozen     & 0.1972     & 0.2093    & 0.1867     & 0.1825     & 0.2198 \\
  & Adaptive   & 0.1994     & 0.2018    & 0.1818     & 0.1939     & 0.2091 \\
  & $\Delta\%$ & $+$1.1     & $-$3.6    & $-$2.6     & $+$6.2     & $-$4.9 \\
\midrule
\multirow{3}{*}{\ddqn{}}
  & Frozen     & 0.2165     & 0.1823    & 0.2195     & 0.2047     & 0.2092 \\
  & Adaptive   & 0.2181     & 0.2133    & 0.2298     & 0.2047     & 0.1696 \\
  & $\Delta\%$ & $+$0.7     & $+$17.0   & $+$4.7     & 0.0        & $-$18.9 \\
\bottomrule
\end{tabular}
\caption{\label{tab:fva}Frozen versus adaptive comparison. Per-episode reward for frozen
  and adaptive modes with relative differences $\Delta\%$ (adaptive
  minus frozen).}
\end{table}

\FloatBarrier

\subsection*{Cross-architecture analysis}
\label{sec:crossarch}

Table~\ref{tab:normalised} presents normalised performance scores
(higher always better; latency and file count inverted and normalised
to [0,~1]).
No architecture is Pareto-dominant, though the five-episode averages
reveal a considerably tighter cluster than single-run figures suggested.
\attentive{} leads on global reward (1.000) and file efficiency (1.000);
\ddqn{} leads on latency (1.000, just 0.2\,ms faster than \attentive{})
and achieves a mid-range pruning score (0.475); \mlpppo{} leads on
pruning ratio (1.000) but trails on both reward and latency.
This complementary profile suggests that architecture selection should
be driven by operational priority: \attentive{} for balanced
optimisation, and \mlpppo{} where partition alignment is the primary concern.
Figure~\ref{fig:rl_vs_baselines} shows the full four-metric profile
across all agent and baseline categories.

\begin{table}[!htbp]
\centering
\begin{tabular}{lccc}
\toprule
\textbf{Metric} & \textbf{\attentive{}} &
  \textbf{\mlpppo{}} & \textbf{\ddqn{}} \\
\midrule
Global reward (higher = better)       & \textbf{1.000} & 0.000          & 0.935 \\
Latency, inverted (lower = better)    & 0.997          & 0.000          & \textbf{1.000} \\
File count, inverted (lower = better) & \textbf{1.000} & 0.000          & 0.507 \\
Pruning ratio (higher = better)       & 0.000          & \textbf{1.000} & 0.475 \\
\bottomrule
\end{tabular}
\caption{\label{tab:normalised}Normalised cross-architecture performance. Each metric is
  normalised to $[0, 1]$ within the three-agent cohort using five-episode
  averaged values; latency and file count are inverted so that higher is
  always better.}
\end{table}

\begin{figure}[!htbp]
\centering
\includegraphics[width=0.92\linewidth]{charts/fig_rl_vs_baselines}
\caption{\label{fig:rl_vs_baselines}\textbf{RL agents vs baselines --- four-metric profile
  on the 1,000-step standard workload.}
  Mean global reward \textbf{(a)}, mean query latency \textbf{(b)},
  mean file count \textbf{(c)}, and mean pruning ratio \textbf{(d)}
  for selected baselines, ablations, and the three RL agents.
  The RL agents (starred) uniformly occupy the top positions on reward
  while achieving competitive latency and file-count metrics.
  PartitionOnly is an extreme outlier on latency (5,307\,ms) and
  file count (1,245).}
\end{figure}

\FloatBarrier

\subsection*{Ablation study}
\label{sec:ablation}

Table~\ref{tab:ablation} compares compaction-only and partition-only
single-agent ablations against all three full multi-agent RL systems.
All ablation agents use the \attentive{} architecture restricted to the
relevant action sub-space.
\texttt{CompactOnly} achieves a reasonable reward (0.1699) and low
latency (344.1\,ms) but zero pruning ratio, confirming compaction alone
cannot adapt query performance to workload characteristics.
\texttt{PartitionOnly} is catastrophic: mean file count 1,245, mean
latency 5,307\,ms, and reward $-0.1765$.
Without compaction, continuous ingestion generates files faster than
any partition repartitioning can reduce them.
This confirms that \textbf{compaction is a prerequisite for effective
partitioning}.
The full multi-agent systems outperform \texttt{CompactOnly} by
$+$25.0--29.5\% and \texttt{PartitionOnly} by $+$220.3--224.7\%.
These gains are non-additive: the meta-controller's joint coordination
provides a synergistic benefit not captured by either sub-agent
independently.

\begin{table}[!htbp]
\centering
\begin{tabular}{lcccc}
\toprule
\textbf{Policy} & \textbf{Reward} & \textbf{Latency (ms)} &
  \textbf{Files} & \textbf{Pruning} \\
\midrule
\texttt{CompactOnly}    & $\phantom{-}0.1699$ & $\phantom{5,}344.1$ &
  $\phantom{1,}48.9$ & 0.000 \\
\texttt{PartitionOnly}  & $-0.1765$           & $5{,}306.9$ &
  $1{,}245.0$ & 0.110 \\
\midrule
\attentive{} & $\mathbf{0.2201}$ & $323.6$ & $54.9$ & $0.191$ \\
\mlpppo{}    & $0.2124$          & $387.8$ & $62.7$ & $0.203$ \\
\ddqn{}      & $0.2196$          & $323.4$ & $58.8$ & $0.197$ \\
\midrule
Gain vs.\ \texttt{CompactOnly}: &&&& \\
\quad \attentive{} & $+$29.5\% &&& \\
\quad \mlpppo{}    & $+$25.0\% &&& \\
\quad \ddqn{}      & $+$29.3\% &&& \\
Gain vs.\ \texttt{PartitionOnly}: &&&& \\
\quad \attentive{} & $+$224.7\% &&& \\
\quad \mlpppo{}    & $+$220.3\% &&& \\
\quad \ddqn{}      & $+$224.4\% &&& \\
\bottomrule
\end{tabular}
\caption{\label{tab:ablation}Ablation study. Compaction-only and partition-only
  single-agent ablations (using the \attentive{} policy restricted to
  the relevant action sub-space) compared against the three full
  multi-agent RL systems. Per cent gain is relative to each ablation baseline.}
\end{table}

\FloatBarrier

\subsection*{Commercial cost analysis}

Table~\ref{tab:cost} presents estimated per-episode cloud costs.
The no-maintenance policy is dominated by query cost (\$42.98 of
\$49.86 total) due to high latency (4,298\,ms mean).
Always-compact policies reduce query cost but incur \$5.00 compaction
overhead, yielding \$15.17 ($-$69.6\%).
The three RL agents achieve the lowest totals: \ddqn{} at \$8.53
($-$82.9\%), \attentive{} at \$8.60 ($-$82.8\%), and \mlpppo{} at
\$9.32 ($-$81.3\%).
RL agents' savings are driven primarily by latency reduction (low query
cost); their compaction costs (\$0.83--\$0.91) are $5\times$ lower
than always-compact policies, reflecting the five-episode-averaged op
counts rather than a single run.

\begin{table}[!htbp]
\centering
\begin{tabular}{lccccc}
\toprule
\textbf{Policy} & \textbf{Query (\$)} & \textbf{Compact (\$)} &
  \textbf{Storage (\$)} & \textbf{Total (\$)} & \textbf{Saving} \\
\midrule
No maintenance   & 42.98 & 0.00 & 6.88 & 49.86 & --- \\
Always compact   &  5.54 & 5.00 & 4.63 & 15.17 & $-$69.6\% \\
Adaptive compact &  6.10 & 4.51 & 4.56 & 15.17 & $-$69.6\% \\
\attentive{}     &  3.24 & 0.91 & 4.45 &  8.60 & $-$82.8\% \\
\mlpppo{}        &  3.88 & 0.91 & 4.53 &  9.32 & $-$81.3\% \\
\ddqn{}          & \textbf{3.23} & \textbf{0.83} & 4.47 &
  \textbf{8.53} & $\mathbf{-82.9\%}$ \\
\bottomrule
\end{tabular}
\caption{\label{tab:cost}Estimated cloud cost breakdown per 1,000-step evaluation
  episode. Cost model: query at \$0.01 per 1,000\,ms; compaction at
  \$0.005 per operation; storage at \$0.001 per MB-step. RL compaction
  counts are averaged over five independent episodes. Savings are
  relative to no-maintenance (\$49.86).}
\end{table}

\begin{figure}[!htbp]
\centering
\includegraphics[width=0.80\linewidth]{charts/fig_cost}
\caption{\label{fig:cost}\textbf{Estimated cloud cost breakdown per evaluation episode.}
  Stacked bars show query cost (red), compaction cost (teal), and
  storage cost (purple) for six policies. Total cost annotations appear
  above each bar; percentage savings relative to no-maintenance appear
  below. No-maintenance is dominated by query cost (\$42.98 of \$49.86).
  Always-compact reduces query cost but incurs \$5.00 compaction
  overhead. The three RL agents achieve the lowest totals
  (\$8.53--\$9.32, savings of 81--83\%) primarily through latency
  reduction rather than minimising compaction operations.}
\end{figure}

\FloatBarrier

% ════════════════════════════════════════════════════════════════════════════
\section*{Discussion}
\label{sec:discussion}
% ════════════════════════════════════════════════════════════════════════════

Our results establish that hierarchical deep RL is an effective and
practical approach to workload-aware lakehouse maintenance.

\textbf{Joint control is essential.}
The $+$224\% reward improvement over \texttt{PartitionOnly} and the
$+$29\% improvement over \texttt{CompactOnly} demonstrate that
compaction and partitioning are complementary operations whose benefits
are only fully realised in combination. A storage layout optimised for
query pruning cannot be sustained without concurrent compaction to
prevent file proliferation.

\textbf{The attention mechanism provides consistent advantages.}
\attentive{} outperforms \mlpppo{} by $+$3.5\% reward and exhibits
lower per-episode variance on the primary metrics (reward CV 5.86\%
versus 6.99\%, latency CV 2.33\% versus 3.24\%). The self-attention encoder's
ability to weight recent observations differentially is advantageous in
the non-stationary workload environment where query profile transitions
demand rapid policy adaptation.

\textbf{Architecture-dependent online adaptation.}
The contrasting adaptive behaviour of \attentive{} ($+$20.3\%) and
\ddqn{} ($-$22.2\%) reveals a fundamental constraint: off-policy
algorithms with experience replay are ill-suited to online adaptation in
a non-stationary environment. When the replay buffer is populated with
transitions from a non-representative initial policy, off-policy gradient
updates reinforce suboptimal behaviours. This finding has practical
implications: frozen weights are robust across architectures, while
online fine-tuning should be restricted to on-policy methods.

\textbf{Rule-based meta-control is sufficient for arbitration.}
The interpretable meta-controller, using simple piecewise-linear urgency
scoring, consistently delegates to the correct specialist. This improves
training stability (specialists receive focused reward signals) and
system transparency (urgency scores can be logged and audited). The
4-step compaction cooldown and 20-step partition cooldown prevent action
thrashing without requiring learned coordination.

\textbf{Limitations.}
The current evaluation is limited to a single table schema and five
query profiles. The cost model uses fixed per-operation prices that may
not reflect cloud provider variability. The RL agents were trained on the
v5 standard workload and evaluated on a compressed version;
generalisation to substantially different workloads (e.g., append-only,
deletion-heavy) has not been tested. Training time and compute
requirements for the three architectures are not reported here and
constitute future work.

\FloatBarrier

% ════════════════════════════════════════════════════════════════════════════
\section*{Methods}
\label{sec:methods}
% ════════════════════════════════════════════════════════════════════════════

\subsection*{The LakeGym v5 environment}

\env{} simulates a production data lakehouse backed by a real Apache
Iceberg table managed via PySpark. Each discrete time step executes:
(i)~ingest a batch of rows into the Iceberg table,
(ii)~the active policy selects one of eight discrete actions,
(iii)~a query is drawn from the current workload profile and executed,
(iv)~three reward signals are computed, and
(v)~an observation vector is constructed.

\textbf{Table schema.}
The table contains eight columns: \texttt{device\_id} (string),
\texttt{timestamp} (timestamp), \texttt{temperature} (double),
\texttt{humidity} (double), \texttt{region} (string),
\texttt{event\_type} (string), \texttt{sensor\_value} (double),
and \texttt{payload} (string, length~100). Each row occupies
approximately 200~bytes; the target file size is 64\,KB, corresponding
to roughly 50~rows per Parquet file.

\textbf{Action space.}
The joint action space consists of eight discrete actions
(Table~\ref{tab:actions}). Compaction actions (\cthirtytwo{},
\csixtyfour{}, \conehundredtwentyeight{}) rewrite small files into
consolidated Parquet files of the specified target size. Partition
actions (\phour{}, \pregion{}, \ptype{}, \removep{}) change the Iceberg
partition specification. \noop{} takes no maintenance action.

\begin{table}[!htbp]
\centering
\begin{tabular}{lcc}
\toprule
\textbf{Action} & \textbf{Type} & \textbf{Cost} \\
\midrule
\texttt{NOOP}                   & ---        & 0.0 \\
\texttt{COMPACT\_32KB}          & Compaction & 0.9 \\
\texttt{COMPACT\_64KB}          & Compaction & 1.0 \\
\texttt{COMPACT\_128KB}         & Compaction & 1.1 \\
\texttt{PARTITION\_HOUR}        & Partition  & 2.0 \\
\texttt{PARTITION\_REGION}      & Partition  & 2.0 \\
\texttt{PARTITION\_EVENT\_TYPE} & Partition  & 2.0 \\
\texttt{REMOVE\_PARTITION}      & Partition  & 0.7 \\
\bottomrule
\end{tabular}
\caption{\label{tab:actions}Action space. All eight discrete actions available to the
  joint agent with their associated normalised cost used in the reward
  function.}
\end{table}

\textbf{Observation space.}
The observation vector at step $t$ includes: file count $f_t$, block
utilisation $u_t$, rows ingested $r_t$, partition pruning ratio $p_t$,
running average pruning ratio $\bar{p}_t$, partition skew $s_t$, query
latency $\ell_t$, query type history fractions (five values), steps
since last compaction $\delta_c$, and steps since last partition change
$\delta_p$.

\textbf{Workload plans.}
The v5 standard workload runs for 1,000 steps divided into five equal
query-profile phases of 199~steps each
(\texttt{time\_heavy}\,$\to$\,\texttt{region\_heavy}\,$\to$\,\texttt{mixed}\,$\to$\,\texttt{type\_heavy}\,$\to$\,\texttt{scan\_heavy}).
Each profile specifies query-type probabilities over five templates
(Table~\ref{tab:profiles}). Ingestion rate follows six phases:
warm-up (40~rows/step, steps 1--100), normal load (75, steps 101--300),
burst (225, steps 301--500), recovery (55, steps 501--700),
normal load 2 (75, steps 701--900), and wind-down (35, steps 901--1,000).
The v5 eval workload is a 2$\times$-compressed version (500 steps) used
for multi-episode consistency and adaptation experiments.
Figure~\ref{fig:workload} illustrates both workload plans.

\begin{table}[!htbp]
\centering
\begin{tabular}{lccccc}
\toprule
\textbf{Profile} & \textbf{Time range} & \textbf{Region filter} &
  \textbf{Sensor lookup} & \textbf{Type filter} & \textbf{Full scan} \\
\midrule
\texttt{time\_heavy}   & \textbf{0.60} & 0.10 & 0.10 & 0.10 & 0.10 \\
\texttt{region\_heavy} & 0.10 & \textbf{0.60} & 0.10 & 0.10 & 0.10 \\
\texttt{mixed}         & 0.20 & 0.20          & 0.20 & 0.20 & 0.20 \\
\texttt{type\_heavy}   & 0.10 & 0.10          & 0.10 & \textbf{0.60} & 0.10 \\
\texttt{scan\_heavy}   & 0.10 & 0.10          & 0.10 & 0.10 & \textbf{0.60} \\
\bottomrule
\end{tabular}
\caption{\label{tab:profiles}Query workload profiles. Each profile specifies the sampling
  probability for each of five query templates. The dominant query type
  is shown in bold for focused profiles.}
\end{table}

\begin{figure}[!ht]
  \centering
  \includegraphics[width=0.85\textwidth]{charts/fig_workload_ingestion}
  \caption{\label{fig:workload}\textbf{Ingestion rate schedule (rows/step).}
  Shaded bands show the min--max range. The burst phase reaches 225~rows/step.}
\end{figure}

\FloatBarrier

\textbf{Pruning matrix.}
File pruning at query time depends on the alignment between the current
partition strategy and the query predicate (Table~\ref{tab:pruning}).
\texttt{TIME\_RANGE} queries prune 92\% of files when
\texttt{HOUR}-partitioned, the highest pruning ratio in the environment.
\texttt{FULL\_SCAN} is irreducible at 0\% under all strategies.

\begin{table}[!htbp]
\centering
\begin{tabular}{lcccc}
\toprule
\textbf{Query type} & \textbf{Hour} & \textbf{Region} &
  \textbf{Event type} & \textbf{None} \\
\midrule
Time range    & \textbf{0.92} & ---           & ---           & --- \\
Sensor lookup & \textbf{0.75} & ---           & ---           & --- \\
Region filter & ---           & \textbf{0.67} & ---           & --- \\
Type filter   & ---           & ---           & \textbf{0.67} & --- \\
Full scan     & ---           & ---           & ---           & --- \\
\bottomrule
\end{tabular}
\caption{\label{tab:pruning}Partition pruning matrix. Fraction of Parquet files skipped
  at query time for each query--partition combination. Dashes indicate
  zero pruning. Only semantically aligned pairs achieve non-zero pruning.}
\end{table}

\begin{figure}[!htbp]
\centering
\includegraphics[width=0.55\linewidth]{charts/fig_pruning_matrix}
\caption{\label{fig:pruning}\textbf{Partition pruning matrix.}
  Fraction of Parquet files skipped at query time for each
  query-type--partition-strategy combination. Cell colour encodes the
  pruning ratio (white~$= 0$, deep navy~$= 1$); dashes indicate zero
  pruning. Time-range queries achieve the highest pruning (92\%) under
  hourly partitioning; full-scan queries are irreducible under all
  strategies.}
\end{figure}

\FloatBarrier

\subsection*{Reward function}

The reward function comprises three independently computed signals.

\textbf{Compact reward} (Eq.~\ref{eq:compact}):
\begin{equation}
r_c = -w_f \tilde{f} + w_u u - w_c \tilde{a}_c + w_\tau \tau(a, r),
\label{eq:compact}
\end{equation}
where $\tilde{f} = \min(f / F_{\max}, 1)$, $u$ is block utilisation,
$\tilde{a}_c$ is the normalised compaction action cost, and
$\tau(a, r)$ is the target-size bonus depending on action $a$ and
ingestion rate $r$. Weights: $w_f = 0.25$, $w_u = 0.20$,
$w_c = 0.10$, $w_\tau = 0.45$.

\textbf{Partition reward} (Eq.~\ref{eq:partition}):
\begin{equation}
r_p = w_p p + w_{\Delta p} \Delta p - w_s \tilde{s}
    - w_{cp} \tilde{a}_p,
\label{eq:partition}
\end{equation}
where $p$ is the current partition pruning ratio,
$\Delta p = \bar{p}_t - \bar{p}_{t-1}$ is the pruning improvement,
$\tilde{s}$ is normalised partition skew, and $\tilde{a}_p$ is the
normalised partition action cost. Weights: $w_p = 0.55$,
$w_{\Delta p} = 0.25$, $w_s = 0.10$, $w_{cp} = 0.10$.

\textbf{Global reward} (Eq.~\ref{eq:global}):
\begin{equation}
r_g = -w_\ell \tilde{\ell} - w_f \tilde{f} + w_p p + w_u u,
\label{eq:global}
\end{equation}
where $\tilde{\ell} = \min(\ell / L_{\max}, 1)$ with
$L_{\max} = 15{,}000$\,ms. Weights: $w_\ell = 0.30$, $w_f = 0.20$,
$w_p = 0.30$, $w_u = 0.20$. The global reward is the primary
optimisation target; compact and partition rewards serve as dedicated
training signals for the respective specialist agents.

\FloatBarrier

\subsection*{Hierarchical multi-agent architecture}

The framework comprises three tiers: the \env{} environment, a
rule-based meta-controller, and two specialist RL agents coordinated
through a shared model registry (Figure~\ref{fig:architecture}).
The three candidate neural architectures plugged into this registry
are compared in Figure~\ref{fig:arch_models}.

\textbf{Meta-controller.}
The meta-controller (implemented in \texttt{meta\_controller.py})
computes two urgency scores at each step:
\begin{equation}
\delta = \begin{cases}
\text{COMPACTION} & \text{if } u_c \geq u_p \geq \theta \\
\text{PARTITION}  & \text{if } u_p > u_c \geq \theta \\
\text{NOOP}       & \text{otherwise,}
\end{cases}
\end{equation}
where $\theta = 0.35$ is the minimum urgency threshold.

Compaction urgency $u_c$ is a piecewise-linear function of file count $f$:
\begin{equation}
u_c(f) = \begin{cases}
0.95 & f \geq 96 \\[2pt]
0.60 + 0.35 \cdot \dfrac{f - 36}{60} & 36 \leq f < 96 \\[6pt]
0.10 + 0.50 \cdot \dfrac{f - 12}{24} & 12 \leq f < 36 \\[2pt]
0    & f < 12,
\end{cases}
\end{equation}
with a $+$0.15 boost when block utilisation falls below 0.55.
A 4-step cooldown prevents repeated compaction.

Partition urgency $u_p$ accumulates from three signals: $+$0.40 if
$\bar{p} < 0.25$, $+$0.30 if the dominant query type fraction exceeds
0.45 (focused workload), and $+$0.10 if the partition layout has not
changed for more than 100 steps. A 20-step cooldown prevents partition
thrashing. Urgency is zeroed if $\bar{p} > 0.50$.

\textbf{Neural architectures.}
All architectures accept input of shape $(\text{batch}, W, F)$ and
return $(\boldsymbol{\pi}, V)$: a softmax action-probability vector
and a scalar state-value estimate
(implemented in \texttt{model\_registry.py}).

\textit{\attentive{}} uses a Transformer-style
encoder\textsuperscript{10}:
input $\to$ Dense(64, ReLU) + trainable positional encoding
$\to$ multi-head self-attention (4 heads, key dimension 64,
dropout 0.1) + residual + LayerNorm
$\to$ FFN (128$\to$64, ReLU, dropout 0.1) + residual + LayerNorm
$\to$ global average pooling $\to$ actor Dense($N_a$, softmax) and
critic Dense(1). Window size $W = 10$, $F = 8$ (compaction) or
$F = 14$ (partition).

\textit{\mlpppo{}} is a feedforward actor-critic:
input $\to$ Flatten $\to$ Dense(256, ReLU) $\to$ Dense(128, ReLU)
$\to$ Dropout(0.1); actor: Dense(64, ReLU) $\to$ Dense($N_a$, softmax);
critic: Dense(64, ReLU) $\to$ Dense(1).
Window size $W = 10$ (compaction) or $W = 20$ (partition).

\textit{\ddqn{}} is a dueling double DQN\textsuperscript{11}:
input $\to$ Flatten $\to$ Dense(256, ReLU) $\to$ BatchNorm $\to$
Dense(128, ReLU) $\to$ BatchNorm; value stream: Dense(64, ReLU) $\to$
$V(s)$; advantage stream: Dense(64, ReLU) $\to$ $A(s,a)$;
$Q(s,a) = V(s) + A(s,a) - \overline{A}(s)$.
Action probabilities are computed as $\text{softmax}(Q)$ for API
uniformity. The model uses an experience replay buffer during training.

Figure~\ref{fig:arch_models} provides a side-by-side comparison of the
three neural architectures sharing a common input--output interface.

\begin{figure}[!htbp]
  \centering
  \begin{subfigure}[t]{0.47\textwidth}
    \centering
    \includegraphics[width=\linewidth]{charts/fig_arch_system}
    \caption{\textbf{Multi-agent system flow.}
    A logic-driven hierarchy where the meta-controller manages specialist
    delegation based on urgency scores.}
    \label{fig:architecture}
  \end{subfigure}
  \hfill
  \begin{subfigure}[t]{0.47\textwidth}
    \centering
    \includegraphics[width=\linewidth]{charts/fig_arch_models}
    \caption{\textbf{Neural candidate comparison.}
    An overview of the different structural paradigms---Attention, MLP,
    and Dueling---evaluated for state-value estimation.}
    \label{fig:arch_models}
  \end{subfigure}
  \caption{System architecture and neural model comparison.}
  \label{fig:both}
\end{figure}

\FloatBarrier

\textbf{Training.}
Agents are trained as multi-agent systems: a single meta-controller
coordinates the compaction specialist (4 actions; $W = 10$, $F = 8$)
and the partition specialist (5 actions; $W = 20$, $F = 14$). Frozen
evaluation fixes all network weights during rollout. Adaptive evaluation
applies online weight updates using PPO for \attentive{} and \mlpppo{},
and online DQN updates for \ddqn{}.

\FloatBarrier

\subsection*{Baseline policies}

We evaluate 14 heuristic compaction policies, 3 workload-aware
heuristics, and 2 ablation policies (Table~\ref{tab:baselines}).
Heuristic policies include always-compact variants (C32, C64, C128),
periodic compaction (every 5 or 10 steps), threshold-triggered
compaction (file count $> 10$ or $> 20$), random compaction, and
adaptive compaction (frequency proportional to file growth rate).
Workload-aware heuristics include \texttt{WorkloadAwareThreshold\_v5},
\texttt{PartitionExploration\_v5}, and \texttt{Random\_Partition\_v5}.
Ablation policies (\texttt{CompactOnly}, \texttt{PartitionOnly}) use
the \attentive{} architecture restricted to the relevant action
sub-space.

\begin{table}[!htbp]
\centering
\renewcommand{\arraystretch}{0.85}
\begin{tabular}{lp{8cm}}
\toprule
\textbf{Policy} & \textbf{Description} \\
\midrule
\texttt{No\_Maintenance}          & Takes no action at any step \\
\texttt{AlwaysCompact}            & Compacts every step with \texttt{64KB} \\
\texttt{AlwaysCompact\_C32}       & Compacts every step with \texttt{32KB} \\
\texttt{AlwaysCompact\_C128}      & Compacts every step with \texttt{128KB} \\
\texttt{Periodic5\_C64}           & Compacts every 5 steps with \texttt{64KB} \\
\texttt{Periodic10\_C64}          & Compacts every 10 steps with \texttt{64KB} \\
\texttt{Periodic10\_C32}          & Compacts every 10 steps with \texttt{32KB} \\
\texttt{Periodic10\_C128}         & Compacts every 10 steps with \texttt{128KB} \\
\texttt{Threshold10\_C32}         & Compaction for file count $> 10$ (\texttt{32KB}) \\
\texttt{Threshold10\_C64}         & Compaction for file count $> 10$ (\texttt{64KB}) \\
\texttt{Threshold10\_C128}        & Compaction for file count $> 10$ (\texttt{128KB}) \\
\texttt{Threshold20\_C64}         & Compaction for file count $> 20$ (\texttt{64KB}) \\
\texttt{Random\_Compact}          & Uniformly samples a compaction action \\
\texttt{AdaptiveCompaction}       & Proportional frequency based on file growth \\
\midrule
\texttt{WorkloadAwareThreshold}   & Dual-aware thresholding and partitioning \\
\texttt{PartitionExploration}     & Sequential cycle of partitioning strategies \\
\texttt{Random\_Partition}        & Uniformly samples a partition action \\
\bottomrule
\end{tabular}
\caption{\label{tab:baselines}\textbf{Heuristic baseline strategies.} This table
details the fourteen rule-based compaction methods and three workload-aware benchmarks
used for comparative analysis. These static policies provide a foundational performance
floor for evaluating the efficacy of our adaptive models.}
\end{table}

\FloatBarrier

\subsection*{Evaluation protocol}

\textbf{Standard evaluation.}
Each policy is evaluated on the 1,000-step v5 standard workload.
For RL agents, frozen weights from the final training checkpoint are
used. Rewards are recalculated post-hoc for heuristic baselines using
the v5 reward function to ensure comparability.

\textbf{Multi-episode evaluation.}
Each RL agent is evaluated over five independent 500-step episodes on
the v5 eval workload. The coefficient of variation (CV) across episodes
quantifies run-to-run reliability.

\textbf{Adaptive evaluation.}
Starting from the frozen checkpoint, agents are evaluated adaptively:
weights are updated at each step using the appropriate online algorithm.
Episode~1 uses the frozen checkpoint; subsequent episodes continue from
the updated weights.

\textbf{Cost model.}
$C = C_q \cdot \sum_t \ell_t + C_c \cdot N_c + C_s \cdot \sum_t
\sigma_t$, where $C_q = \$0.00001$/ms, $C_c = \$0.005$ per compaction
operation, $C_s = \$0.000001$/KB-step, $\ell_t$ is per-step latency,
$N_c$ is total compaction operations, and $\sigma_t$ is total table
size in KB at step $t$.

\textbf{Statistical reporting.}
Descriptive statistics are reported as mean $\pm$ s.d.\ or mean with CV
for multi-episode comparisons. No formal hypothesis testing is performed;
comparisons are reported as absolute and relative differences in mean
metric values. The sample size for the primary 1,000-step RL evaluation is
$n = 5{,}000$ total step-level observations per agent (five episodes
of 1,000 steps each), averaged to a representative 1,000-step episode.
For multi-episode comparisons, $n = 5$ episode-level means per agent.

\FloatBarrier

% ════════════════════════════════════════════════════════════════════════════
\section*{References}
% ════════════════════════════════════════════════════════════════════════════

\begin{enumerate}

\item Blue, R., Weeks, D., Reid, J. \& Steinbach, C. Iceberg: A fast and
  expressive table format for large, slow tables. In \textit{Proc. ACM
  SIGMOD} 3434--3441 (ACM, 2020).

\item Armbrust, M. \textit{et al.} Delta Lake: High-performance ACID table
  storage over cloud object stores. \textit{Proc. VLDB Endow.}
  \textbf{13}, 3411--3424 (2020).

\item Chandar, V., Rao, B. S. \& Agarwal, N. Apache Hudi: The data lake
  platform. \textit{arXiv} preprint arXiv:2301.09517 (2023).

\item Kraska, T. \textit{et al.} The case for learned index structures. In
  \textit{Proc. ACM SIGMOD} 489--504 (ACM, 2018).

\item Agrawal, S. \textit{et al.} Database tuning advisor for Microsoft SQL
  Server 2005. In \textit{Proc. VLDB} 1110--1121 (2004).

\item Marcus, R. \& Papaemmanouil, O. Re-Join: How the machine learned to
  join. In \textit{Workshop on ML for Progressively Learned Databases
  (NIPS DBLearn)} (2018).

\item Trummer, I. \textit{et al.} SkinnerDB: Regret-bounded query evaluation
  via reinforcement learning. In \textit{Proc. ACM SIGMOD} 1153--1170
  (ACM, 2019).

\item Van Aken, D., Pavlo, A., Gordon, G. J. \& Zhang, B. Automatic
  database management system tuning through large-scale machine learning.
  In \textit{Proc. ACM SIGMOD} 1009--1024 (ACM, 2017).

\item Zhang, J. \textit{et al.} An end-to-end automatic cloud database
  tuning system using deep reinforcement learning. In \textit{Proc. ACM
  SIGMOD} 415--432 (ACM, 2019).

\item Vaswani, A. \textit{et al.} Attention is all you need. In
  \textit{Advances in Neural Information Processing Systems (NeurIPS)}
  vol.~30 (2017).

\item Wang, Z. \textit{et al.} Dueling network architectures for deep
  reinforcement learning. In \textit{Proc. ICML} 1995--2003 (2016).

\item Dayan, N., Athanassoulis, M. \& Idreos, S. Dostoevsky: Better
  space-time trade-offs for LSM-tree based key-value stores via adaptive
  removal of superfluous merging. In \textit{Proc. ACM SIGMOD} 505--520
  (ACM, 2018).

\item Schulman, J. \textit{et al.} Proximal policy optimization algorithms.
  \textit{arXiv} preprint arXiv:1707.06347 (2017).

\item Camacho-Rodr\'{\i}guez, J. \textit{et al.} LST-Bench: Benchmarking
  Log-Structured Tables in the Cloud. In \textit{Proc. ACM SIGMOD}
  1--14 (ACM, 2024).

\item Armbrust, M. \textit{et al.} Lakehouse: A New Generation of Open
  Platforms that Unify Data Warehousing and Advanced Analytics. In
  \textit{Proc. CIDR} (2021).

\item Yang, Z. \textit{et al.} Qd-tree: Learning Data Layouts for Big
  Data Analytics. In \textit{Proc. ACM SIGMOD} 170--184 (ACM, 2020).

\item Kraska, T. \textit{et al.} SageDB: A Learned Database System. In
  \textit{Proc. CIDR} (2019).

\item Pavlo, A. \textit{et al.} Self-Driving Database Management Systems.
  In \textit{Proc. CIDR} (2017).

\item Marcus, R. \textit{et al.} Neo: A Learned Query Optimizer.
  \textit{Proc. VLDB Endow.} \textbf{12}, 1705--1718 (2019).

\item Li, G. \textit{et al.} QTune: A Query-Aware Database Tuning System
  with Deep Reinforcement Learning. \textit{Proc. VLDB Endow.}
  \textbf{12}, 2118--2130 (2019).

\item Kanellis, K. \textit{et al.} LlamaTune: Sample-Efficient DBMS
  Configuration Tuning. \textit{Proc. VLDB Endow.} \textbf{15},
  2726--2738 (2022).

\item Cai, B. \textit{et al.} HUNTER: An Online Cloud Database Hybrid
  Tuning System for Personalized Requirements. In \textit{Proc. ACM
  SIGMOD} 625--638 (ACM, 2022).

\item Yang, Z. \textit{et al.} HWTune: A Hardware-Aware Database Tuning
  System with Deep Reinforcement Learning. \textit{IEEE Trans. Knowl.
  Data Eng.} \textbf{36}, 2284--2297 (2023).

\item Feng, Y. \textit{et al.} CTuner: Automatic NoSQL Database Tuning
  with Causal Reinforcement Learning. In \textit{Proc. ACM SIGKDD}
  1--12 (ACM, 2024).

\item Liu, Y. \textit{et al.} MLETune: Streamlining Database Knob Tuning
  via Multi-LLMs Experts Guided Deep Reinforcement Learning.
  \textit{IEEE Access} \textbf{12}, 1--15 (2024).

\item Zhang, K., Yang, Z. \& Ba\c{s}ar, T. Multi-Agent Reinforcement
  Learning: A Selective Overview of Theories and Algorithms.
  \textit{J. Mach. Learn. Res.} \textbf{22}, 1--85 (2021).

\item Vezhnevets, A. S. \textit{et al.} FeUdal Networks for Hierarchical
  Reinforcement Learning. In \textit{Proc. ICML} 3540--3549 (2017).

\item Wang, Y. \textit{et al.} A highly efficient distributed database
  knobs tuning method based on multi-agent deep reinforcement learning.
  \textit{J. Supercomput.} (2025).

\end{enumerate}

% ════════════════════════════════════════════════════════════════════════════
\section*{Data availability}
% ════════════════════════════════════════════════════════════════════════════
The datasets generated and analysed during the current study,
including all experimental metrics collected across workload profiles,
are available in the LakePilot repository at
\url{https://github.com/OmarEssameldinMousa/LakePilot-Research-prototype}.
The simulation environment (\env{}) used to generate all data,
together with model weights, meta-controller implementations, and
the model registry, is included in the same repository.
No pre-existing external datasets were used in this study; all data
was synthetically generated by the \env{} environment as described
in the Methods section.
% ════════════════════════════════════════════════════════════════════════════
\section*{Acknowledgements}
% ════════════════════════════════════════════════════════════════════════════

The authors developed the \env{} simulation environment utilising the
open-source Apache Iceberg and Apache Spark frameworks.
All figures were generated using Matplotlib
(v3.8, \url{https://matplotlib.org}).


\section*{Funding}
This research did not receive any specific grant from funding agencies
in the public, commercial, or not-for-profit sectors.

% ════════════════════════════════════════════════════════════════════════════
\section*{Author contributions statement}
% ════════════════════════════════════════════════════════════════════════════

O.E.\ conceived the idea, implemented the \env{} environment,
developed the lakehouse solution, and    designed the whole experiments.
O.K.\ implemented and trained the AI models.
O.S.\ assisted in baseline analysis and overall project execution.
A.Z.B.\ supervised the research.
All authors reviewed and approved the final manuscript.

% ════════════════════════════════════════════════════════════════════════════
\section*{Additional information}
% ════════════════════════════════════════════════════════════════════════════

\textbf{Competing interests.} The authors declare no competing interests.

\end{document}